% A Gaussian Evolutionary Strategy for Binary Weight Neural Network Optimization
% EvoStar 2027 (EvoApplications track) submission.
%
% Companion paper to the Bernoulli ES paper (../bernoulli-paper): same problem,
% same evaluation protocol, but the search distribution is Gaussian over
% continuous surrogate parameters, and perturbed parameters are quantized with
% the sign function before every fitness evaluation.
%
% Build: pdflatex total && bibtex total && pdflatex total && pdflatex total
%
% TODO before submission:
%   - author names, affiliations, ORCIDs (search "TODO")
%   - acknowledgments / funding
%   - code/data URL for the reproducibility statement
%   - all numbers, figures, and tables (placeholders below)
\documentclass[runningheads]{llncs}

\usepackage[T1]{fontenc}
\usepackage{graphicx}
\usepackage{amsmath}
\usepackage{amssymb}
\usepackage{booktabs}
\graphicspath{{figures/}}

\begin{document}

\title{A Gaussian Evolutionary Strategy for Binary Weight Neural Network Optimization}

% TODO: real author names, ORCIDs and institutes
\author{}
\authorrunning{}
\institute{}

\maketitle

% TODO: draft abstract.
% Key points to hit:
%   - binary weight networks: same setting as the Bernoulli paper
%   - search distribution is Gaussian over continuous surrogate parameters
%     (latent weights / inertia interpretation), NOT over the binary vectors
%   - perturbed parameters are quantized with sign(), so fitness is always
%     evaluated on a valid binary network
%   - REINFORCE score reduces to eps/sigma (pre-quantization sampling noise)
%   - because the search parameters are plain continuous variables updated by
%     a scalar gradient estimate, standard optimizers such as Adam apply in
%     exactly the same way as in the straight-through-estimator pipeline
%   - MNIST results: population scaling, comparison to Adam-STE baseline
\begin{abstract}
TODO
\keywords{Binary Weight Neural Networks \and Evolutionary Strategies \and Quantization}
\end{abstract}

\section{Introduction}
\label{sec:intro}

% Structure (mirroring the Bernoulli paper, with the positioning flipped):
%   para 1: binary weight networks, memory motivation, training via
%           straight-through estimator (cite as in the Bernoulli paper).
%   para 2: interpretations of surrogate parameters: inertia
%           (Helwegen et al., latent weights) vs. Bernoulli expectations
%           (Meng et al., Bayesian learning rule); our Bernoulli paper took
%           the expectation interpretation literally and optimized a
%           distribution over binary vectors directly.
%   para 3: this work takes the inertia interpretation literally instead:
%           maintain continuous surrogate parameters theta, perturb them with
%           isotropic Gaussian noise, quantize sign(theta + sigma eps) for
%           every fitness evaluation. This is an evolutionary strategy whose
%           search distribution lives in continuous parameter space, so the
%           latent parameters can be optimized with Adam exactly like the
%           surrogate parameters of the STE pipeline -- connecting ES and STE
%           practice. Positioning w.r.t. OpenAI ES (Salimans et al.) and
%           quantized / quantization-aware ES literature.

TODO

\section{Binary Weight Neural Networks}
\label{sec:BWNN}

% Same as the Bernoulli paper: sign-mapped binary weights, per-layer scale
% sqrt(2)/sqrt(n_in) replacing batch normalization for ReLU networks, fitness
% = negative categorical cross-entropy on random mini-batches.
% Reuse that text with light edits.

TODO

\section{Gaussian Evolutionary Strategy}
\label{sec:GES}

% Derivation sketch (to be written out in full):
%
% Search parameters: continuous vector theta in R^D (the surrogate/latent
% parameters). Sampling: x_i = theta + sigma*eps_i, eps_i ~ N(0, I_D).
% Quantization: the network weights of individual i are w_i = sign(x_i), so
% every fitness evaluation is on a valid binary network and the fitness
% function f: {-1,+1}^D -> R is exactly the same black box as in the Bernoulli
% paper.
%
% Objective: J(theta) = E_{eps}[ f(sign(theta + sigma eps)) ].
% J is smooth in theta through the sampling noise.
%
% REINFORCE gradient:
%   nabla_theta log N(x_i; theta, sigma^2 I) = (x_i - theta)/sigma^2 = eps_i/sigma,
% so the population-summed update is
%   theta <- theta + eta/(P sigma) * sum_i u_i eps_i
% with shaped fitness u_i. Note the score does not depend on the network at
% all: only the pre-quantization sampling noise enters the estimator.
%
% Points to work out in the text:
%   - sign(0) convention: locked decision -- the search parameters are
%     initialized from a zero-mean Gaussian N(0, init_std^2), so exact zeros
%     (where the sign map is undefined) have measure zero; bits are quantized
%     as x > 0. init_std is a new hyperparameter relative to the Bernoulli
%     paper (default 1.0; with init_std = sigma the initial perturbation flip
%     probability is ~50%).
%   - sigma: LOCKED as a shared fixed scale (no per-coordinate sigma_j);
%     commitment/saturation is then governed by |theta_j|/sigma growth
%   - antithetic sampling is implemented (eps paired with -eps) but default
%     off and not studied in this paper
%   - optimizer preprocessing: because the update is a plain additive
%     continuous gradient, Adam/momentum apply unchanged -- the algorithm is
%     "STE with the gradient replaced by a REINFORCE estimate", which is the
%     headline connection
%   - memory/cost parity w.r.t. the Bernoulli ES: the Bernoulli sampler also
%     generates a (P, D) float tensor per generation (it just never keeps it),
%     so sampling cost is the same; the only Gaussian-specific cost is
%     re-drawing the noise at update time, made exact by replaying the same
%     per-generation seeded stream chunk-by-chunk (no (P, D) float is ever
%     materialized whole) -- worth a remark, not a section
%   - relation to the Bernoulli ES: the induced bit distribution is a product
%     of Bernoullis with rates p_j = Phi(theta_j/sigma) (erf), i.e. the same
%     Bernoulli family under a monotone logit reparameterisation -- the two
%     algorithms differ in parameterisation, score function (eps/sigma vs the
%     bounded x - s), and convergence dynamics (commitment appears as
%     |theta_j| >> sigma growth instead of logit saturation); diagnostics
%     (entropy, committed fraction, per-layer histograms over p_j) are
%     computed identically to the Bernoulli paper, giving direct figure parity

TODO

\section{Methods}
\label{sec:methods}

% Same protocol as the Bernoulli paper (MNIST 54k/6k/10k split, standardised
% inputs, fully connected ReLU MLP with pre-activation scaling, no biases,
% fitness = negative CE on random 4096-sample mini-batches, mode-network
% validation evals every 10 iterations, test at end of training).
%
% TODO: decide and record the experimental grid:
%   - sigma (shared fixed; maybe a small sigma sweep)
%   - population sizes, learning rates, Adam vs. plain update for the ES
%   - antithetic on/off
%   - shaping: rank / standardize / none (pick one as default)
%   - architecture, seeds, iterations (match Bernoulli paper: h256x2,
%     5 seeds, 2000 it pop sweep, 5000 it traces)
%   - Adam-STE baseline: Larq, same arch, Adam 1e-3, batch 512, 100 epochs

TODO

\section{Results}
\label{sec:results}

% Planned figures/tables (placeholders):
%   - fig_pop_scaling: val/test accuracy + entropy vs. population size
%   - fig_val_traces: per-seed validation traces
%   - theta/|theta| histograms per layer over training: saturation as
%     |theta_j|/sigma growth (the Gaussian analogue of the Bernoulli
%     expectation histograms)
%   - table: Gaussian ES vs. Adam-STE vs. (optionally) Bernoulli ES numbers

TODO

\section{Discussion}
\label{sec:discussion}

TODO

\section{Conclusion}
\label{sec:conclusion}

TODO

\begin{credits}
% TODO: funding/acknowledgments or remove

\subsubsection{\discintname}
The authors have no competing interests to declare that are relevant to the
content of this article.
\end{credits}

\bibliographystyle{splncs04}
\bibliography{refs}

\end{document}
